% ============================================================
%  01 · Grundlagen zur Anwendung von Zufallsvariablen
%      Quelle: 01_Grundlagen.pdf
% ============================================================

\bereich[cMath]{Mathematische Grundregeln}
\noindent\footnotesize\setlength{\tabcolsep}{6pt}%
\begin{tabular}{@{}l|l@{}}
  $\sqrt[n]{x}=x^{1/n}$ &
  $\dfrac{1}{x^n}=x^{-n}$ \\[0.3em]
  \multicolumn{2}{@{}l}{Kettenregel (Klammerpotenz):\quad
    $\bigl[u(x)^{n}\bigr]' =
      \underbrace{n\cdot u(x)^{n-1}}_{\text{\scriptsize äuß.\ Abl.}}
      \cdot
      \underbrace{u'(x)}_{\text{\scriptsize inn.\ Abl.}}$} \\[0.3em]
  \multicolumn{2}{@{}l}{\rule{0pt}{1.1em}%
    \textbf{Gerade $y(x)$ aus Graph ablesen:}\quad
    $y(x)=m\,x+b$} \\[0.2em]
  $b$: Schnittpunkt mit $y$-Achse (Wert bei $x{=}0$) &
  $m=\dfrac{\Delta y}{\Delta x}=\dfrac{y_2-y_1}{x_2-x_1}$
    \;(zwei Punkte) \\[0.4em]
  \multicolumn{2}{@{}l}{Punkt-Steigungs-Form:\quad
    $y(x)=y_1+m\,(x-x_1)$\qquad
    {\scriptsize($m{>}0$ steigend, $m{<}0$ fallend, $m{=}0$ waagrecht)}} \\
\end{tabular}\par
\bereichende

\bereich[cMath]{Typische Funktionen im Graphen erkennen}
\begin{formelreihe}
  \parbox[t]{\linewidth}{%
    {\scriptsize\scshape\color{gray!65!black} Dreieck}\\[0.03em]%
    $\displaystyle f(\tau)=A\!\left(1-\dfrac{|\tau-\tau_0|}{b}\right)$\\[0.1em]%
    {\tiny\itshape für $|\tau{-}\tau_0|\le b$, sonst $0$. Höhe $A$, Zentrum
      $\tau_0$, Halbbreite $b$, Flankensteig.\ $A/b$. Erkennung: gerade
      Flanken, Knick an der Spitze. Entsteht bei Faltung zweier gleicher
      Rechtecke.}%
  } &
  \parbox[t]{\linewidth}{%
    {\scriptsize\scshape\color{gray!65!black} Rechteck}\\[0.03em]%
    $\displaystyle f(\tau)=A \;\text{für}\; a\le\tau\le c,\;\text{sonst }0$\\[0.1em]%
    {\tiny\itshape Höhe $A$, Grenzen $a,c$, Breite $c-a$.}%
  } &
  \parbox[t]{\linewidth}{%
    {\scriptsize\scshape\color{gray!65!black} Rampe / Gerade}\\[0.03em]%
    $\displaystyle m=\dfrac{y_2-y_1}{x_2-x_1},\;\; f(x)=y_1+m(x{-}x_1)$\\[0.1em]%
    {\tiny\itshape im gültigen Bereich.}%
  } \\[0.3em]
  \parbox[t]{\linewidth}{%
    {\scriptsize\scshape\color{gray!65!black} Exponentialabfall}\\[0.03em]%
    $\displaystyle f(\tau)=A\,e^{-|\tau|/T}$\\[0.1em]%
    {\tiny\itshape Höhe $A$ bei $\tau{=}0$, Zeitkonstante $T$ (Wert $A/e$ bei
      $|\tau|{=}T$). Typische AKF von tiefpassgefiltertem Rauschen.}%
  } &
  \parbox[t]{\linewidth}{%
    {\scriptsize\scshape\color{gray!65!black} Dirac-Impuls}\\[0.03em]%
    $\displaystyle f(\tau)=A\,\delta(\tau-\tau_0)$\\[0.1em]%
    {\tiny\itshape Pfeil der Gewichtung $A$ bei $\tau_0$. Weißes Rauschen:
      $\varphi(\tau)=N_0\,\delta(\tau)$.}%
  } &
  \parbox[t]{\linewidth}{%
    {\scriptsize\scshape\color{gray!65!black} Cosinus}\\[0.03em]%
    $\displaystyle f(\tau)=A\cos(\omega_0\tau),\;\; \omega_0=\dfrac{2\pi}{T}$\\[0.1em]%
    {\tiny\itshape $\omega_0$ aus dem Abstand zweier Maxima. AKF eines
      harmonischen Prozesses.}%
  } \\[0.3em]
  \multicolumn{3}{@{}p{\linewidth}}{%
    \infoblock{\textbf{Faustregeln:}\;
      Gerade Flanken $\to$ Betragsfunktion (Dreieck); gekrümmte Flanken $\to$
      Quadrat (Parabel).\quad
      Jede AKF ist gerade ($\varphi(\tau)=\varphi(-\tau)$) mit Maximum bei
      $\tau=0$ — Plausibilitätscheck beim Ablesen.}%
  } \\
\end{formelreihe}
\bereichende

\bereich[cGrundl]{Wahrscheinlichkeiten}
\begin{formelreihe}
  \infoblock{$\Omega$: Ergebnismenge,\; $E\subseteq\Omega$: Ereignis.\\
             Würfel: $\Omega=\{1,\dots,6\}$, gerade $=\{2,4,6\}$} &
  \formelblock{Wahrscheinlichkeit}{P(E)=\dfrac{|H|}{|\Omega|}=\dfrac{N_{Elem.Ereignismenge}}{N_{Elem.Ergebnismenge}}} &
  \beispielblock{Würfel}{P(\text{gerade})=\tfrac{3}{6}=\tfrac{1}{2},\;
    P(X{=}3)=\tfrac{1}{6}} \\[0.3em]
  \formelblock{Komplement}{P(\bar{E})=1-P(E)} &
  \formelblock{Additivität (disjunkt)}{P(E\cup F)=P(E)+P(F)} &
  \beispielblock{Komplement}{P(X{\le}4)=1-P(X{>}4)=1-\tfrac{1}{3}=\tfrac{2}{3}} \\
\end{formelreihe}
\bereichende

\bereich[cGrundl]{Bedingte Wahrscheinlichkeit \& Bayes}
\begin{formelreihe}
  \formelblock{Bedingte W'keit}{P(B\mid A)=\dfrac{P(B\cap A)}{P(A)}} &
  \formelblock{Satz von Bayes}{P(A\mid B)=\dfrac{P(B\mid A)\,P(A)}{P(B)}} &
  \beispielblock{$B{=}\{5,6\}$, $A{=}\{2,4,6\}$}{%
    P(B{\mid}A)=\dfrac{P(\{6\})}{P(\{2,4,6\})}
    =\dfrac{1/6}{3/6}=\dfrac{1}{3}} \\[0.3em]
  \formelblock{Unabhängigkeit}{P(A\cap B)=P(A)\cdot P(B)} &
  \formelblock{Folge bei Unabh.}{P(B\mid A)=P(B)} &
  \beispielblock{$B{=}\{5,6\}$, $A{=}\{2,4,6\}$}{P(A{\mid}B)=\dfrac{\frac13\cdot\frac12}{\frac13}=\dfrac{1}{2}} \\
\end{formelreihe}
\bereichende

\bereich[cGrundl]{Zufallsvariablen}
\begin{formelreihe}
  \infoblock{$X:\Omega\to\mathbb{R}$ ordnet jedem Ergebnis eine reelle Zahl zu.\\
             Diskret: endlich/abzählbar viele Werte.\\
             Kontinu.: $P(X{=}x)=0\;\forall\,x$, z.\,B.\ \\Glücksrad $X\in[0,2\pi]$} &
  \formelblock{Verteilungsfunktion}{F_X(x)=P(X\le x)} &
  \beispielblock{Würfel}{X\in\{1,\dots,6\},\quad P(X{=}x)=\tfrac{1}{6}} \\[0.3em]
  \formelblock{Intervall-W'keit}{P(a{<}X{\le}b)=F_X(b)-F_X(a)} &
  \formelblock{Überschreitung}{P(X{>}x)=1-F_X(x)} &
  \\
\end{formelreihe}
\bereichende

\bereich[cGrundl]{Wahrscheinlichkeitsdichte (PDF) \& Verteilung (CDF)}
% Spalte 1 = Dichte (PDF) · Spalte 2 = Verteilung (CDF) · Spalte 3 = Beispiele
\begin{formelreihe}
  \formelblock{PDF: Normierung}{\int_{-\infty}^{\infty}\!f_X(x)\,\mathrm{d}x=1,
    \quad f_X(x)\ge0} &
  \formelblock{CDF: Definition}{F_X(x)=\int_{-\infty}^{x} f_X(\xi)\,\mathrm{d}\xi} &
  \beispielblock{Gleichvert.\ $[0,2\pi]$, $f_X{=}\tfrac{1}{2\pi}$}{%
    P\!\left(\tfrac{\pi}{2}{<}X{<}\tfrac{3\pi}{4}\right)
    =F_X\!\!\left(\tfrac{3\pi}{4}\right)\!-\!F_X\!\!\left(\tfrac{\pi}{2}\right)
    =\tfrac{1}{8}} \\[0.3em]
  \formelblock{Diskrete PDF}{f_X(x)=\sum_i P(X{=}x_i)\cdot\delta(x-x_i)} &
  \formelblock{Ränder \& W'keiten}{\begin{aligned}
    F_X(-\infty)&=0,\quad F_X(\infty)=1 \\
    P(X{\le}x)&=F_X(x)\\
    P(X{>}x)&=1-F_X(x)\\
    P(x_1{<}X{\le}x_2)&=F_X(x_2)-F_X(x_1)
  \end{aligned}} &
  \\[0.3em]
  \formelblock{Gaußdichte}{f_X(x)=\dfrac{1}{\sqrt{2\pi}\,\sigma_X}
    \exp\!\!\left(-\dfrac{1}{2}\!\left(\dfrac{x-\mu_X}{\sigma_X}\right)^{\!2}\right)} &
  \formelblock{Gauß-CDF \& erf}{F_X(x)=\tfrac{1}{2}\!\left(1+\mathrm{erf}\!\left(
    \dfrac{x-\mu_X}{\sqrt{2}\,\sigma_X}\right)\right)\\[0.2em]
    \mathrm{erf}(x)=\dfrac{2}{\sqrt{\pi}}\int_0^x e^{-\xi^2}\,\mathrm{d}\xi} &
  \\
\end{formelreihe}
\bereichende

\bereich[cGrundl]{REZEPT · Gültige WDF / Normierung (+ Mittelwertfreiheit)}
\begin{rezept}
  \rs{Fläche $=1$ ansetzen: $\int f_X\,\mathrm dx=1$.\;
      Rechteck: Breite${\cdot}$Höhe; Dreieck: $\tfrac12{\cdot}$Breite${\cdot}$Höhe;
      Diracs: Summe der Gewichte $\sum w_i=1$.}
  \rs{Unbekannte $>1$? Zweite Gleichung aus Zusatzbedingung — meist
      \emph{Mittelwertfreiheit} $\int x\,f_X\,\mathrm dx=0$ oder gegebene Varianz.}
  \rs{Lineares Gleichungssystem lösen.}
\end{rezept}
\miniexample{Normierung (a) \;/\; Normierung + Mittelwertfreiheit (b)}{%
  (a) Rechteck Höhe $f_0$ auf $[0,5]$ + Dirac $\tfrac14$ bei $x{=}2$:\;
  $5f_0+\tfrac14=1\Rightarrow f_0=\tfrac{3}{20}$.\hfill
  (b) Rechtecke $f_1$ auf $[-4,0]$, $f_2$ auf $[0,2]$:\;
  $4f_1{+}2f_2{=}1$ und $\mu_X{=}0:\;-8f_1+2f_2=0$
  $\Rightarrow f_2{=}4f_1\Rightarrow f_1{=}\tfrac1{12},\;f_2{=}\tfrac13$.}
\bereichende

\bereich[cGrundl]{Erwartungswert, Momente \& Varianz}
\begin{formelreihe}
  \formelblock{Erwartungs-/Mittelwert (stetig)}{E\{X\}=\mu_X
    =\int_{-\infty}^{\infty} x\,f_X(x)\,\mathrm{d}x} &
  \formelblock{Erwartungswert (diskret)}{E\{X\}=\mu_X
    =\sum_i x_i\,P(X{=}x_i)} &
  \beispielblock{Würfelmittel}{E\{X\}=\sum_{k=1}^{6}k\cdot\tfrac{1}{6}
    =\tfrac{21}{6}=3{.}5} \\[0.3em]
  \parbox[t]{\linewidth}{%
    {\scriptsize\scshape\color{gray!65!black} EW diskret via Dirac}\\[0.03em]%
    {\scriptsize$\displaystyle E[X]=\int_{-\infty}^{\infty}\!x\cdot\sum_i w_i\cdot\delta(x{-}x_i)\,\mathrm{d}x=\left[w_i\cdot x_i\right]_{-\infty}^{\infty}$}%
  } &
  \infoblock{$w_i=P(X{=}x_i)$\\[0.1em]
    Die Ausblendeigenschaft ersetzt x durch $x_i$, \\der Rest fällt weg.} &
  \\[0.3em]
  \formelblock{$n$-tes Moment}{m_X^{(n)}=E\{X^n\}
    =\int_{-\infty}^{\infty}x^n f_X(x)\,\mathrm{d}x} &
  \formelblock{$n$-tes zentrales Moment}{\mu_X^{(n)}
    =E\!\left\{\!\left(X-m_X^{(1)}\right)^{\!n}\right\}} &
  \\[0.3em]
  \parbox[t]{\linewidth}{%
    {\scriptsize\scshape\color{gray!65!black} Varianz (stetig)}\\[0.03em]%
    $\displaystyle\sigma_X^2=\int_{-\infty}^{+\infty}\!(x{-}\mu_X)^2 f_X(x)\,\mathrm{d}x$\\[0.25em]%
    {\scriptsize\scshape\color{gray!65!black} Verschiebungssatz (Steiner)}\\[0.03em]%
    $\displaystyle\sigma_X^2=E\{X^2\}-\mu_X^2$\\[0.25em]%
    {\scriptsize\scshape\color{gray!65!black} Standardabweichung}\\[0.03em]%
    $\displaystyle\sigma_X=\sqrt{E\{(X-\mu_X)^2\}}$\\[0.25em]%
    {\scriptsize\scshape\color{gray!65!black} Gleichverteilung ($f_X$ konstant)}\\[0.03em]%
    $\displaystyle\mu_X=\dfrac{x_{\min}+x_{\max}}{2},\quad\sigma_X^2=\dfrac{(x_{\max}-x_{\min})^2}{12}$%
  } &
  \parbox[t]{\linewidth}{%
    {\scriptsize\scshape\color{gray!65!black} Varianz (diskret)}\\[0.03em]%
    $\displaystyle\sigma_X^2=\sum_i(x_i{-}\mu_X)^2 P(X{=}x_i)$\\[0.25em]%
    {\scriptsize\scshape\color{gray!65!black} Varianz diskret via Dirac}\\[0.03em]%
    {\scriptsize$\displaystyle\sigma_X^2=\left[w_i\cdot x_i^2\right]_{-\infty}^{\infty}-\left(E[X]\right)^2$}\\[0.15em]%
    {\tiny\itshape$w_i=P(X{=}x_i)$}%
  } &
  \parbox[t]{\linewidth}{%
    \beispielblock{Würfelvarianz}{%
      \sigma_X^2=\tfrac{1}{6}\sum_{k=1}^{6}(k-3{.}5)^2
      =\tfrac{17.5}{6}\approx2{.}92}\\[0.2em]%
    \beispielblock{Glücksrad $[0,2\pi]$, $f_X{=}\tfrac{1}{2\pi}$}{%
      \mu_X=\int_0^{2\pi}\!x\cdot\tfrac{1}{2\pi}\,\mathrm{d}x
      =\tfrac{1}{2\pi}\!\left[\tfrac{x^2}{2}\right]_0^{2\pi}=\pi\\[0.2em]
      \sigma_X^2=\int_0^{2\pi}\!(x{-}\pi)^2\cdot\tfrac{1}{2\pi}\,\mathrm{d}x
      =\dfrac{\pi^2}{3}}%
  } \\
\end{formelreihe}
\bereichende

\bereich[cGrundl]{REZEPT · Erwartungswert \& Varianz}
\begin{formelreihe}
  \parbox[t]{\linewidth}{%
    {\scriptsize\scshape\color{gray!65!black} Diskret}\\[0.05em]
    {\footnotesize $\mu=\sum_i x_i w_i$,\quad
      $\sigma^2=\sum_i x_i^2 w_i-\mu^2$}} &
  \parbox[t]{\linewidth}{%
    {\scriptsize\scshape\color{gray!65!black} Stetig (bereichsweise)}\\[0.05em]
    {\footnotesize $\mu=\int x f\,\mathrm dx$,\quad
      $\sigma^2=\int x^2 f\,\mathrm dx-\mu^2$}} &
  \parbox[t]{\linewidth}{%
    {\scriptsize\scshape\color{gray!65!black} Gleichverteilung $[a,b]$}\\[0.05em]
    {\footnotesize $\mu=\tfrac{a+b}{2}$,\quad $\sigma^2=\tfrac{(b-a)^2}{12}$}} \\
\end{formelreihe}
\miniexample{Diskret: $X\in\{0,2,4,6\}$, je $w_i{=}\tfrac14$}{%
  $\mu_X=\tfrac14(0{+}2{+}4{+}6)=3$;\quad
  $E\{X^2\}=\tfrac14(0{+}4{+}16{+}36)=14$;\quad
  $\sigma_X^2=14-3^2=5$ \;(Verschiebungssatz).}
\bereichende

\bereich[cGrundl]{Transformation \& Lineare Abbildungen}
\begin{formelreihe}
  \formelblock{Lineare Abbildung}{Y=aX+b} &
  \formelblock{Mittelwert}{\mu_Y=a\,\mu_X+b} &
  \formelblock{Varianz}{\sigma_Y^2=a^2\,\sigma_X^2} \\[0.3em]
  \formelblock{Dichte-Transformation}{f_Y(y)=\dfrac{f_X(x)}{|y'(x)|}} &
  \formelblock{CDF-Transformation}{F_Y(y)=F_X(g^{-1}(y))\\[0.1em]
    {\scriptsize(g^{-1}\text{: Umkehrfunktion})}} &
  \infoblock{Bereich von $f_Y$: $x$-Grenzen durch\\
    Einsetzen in $y=g(x)$ umrechnen!} \\
\end{formelreihe}
\bereichende

\bereich[cGrundl]{REZEPT · Transformation über Kennlinie $f_Y(y)=\dfrac{f_X(x)}{|g'(x)|}$}
\begin{rezept}
  \rs{Kennlinie stückweise linear $\Rightarrow$ $g'(x)$ = Steigung je Abschnitt
      (konstant).}
  \rs{Pro $y$-Bereich den zugehörigen $x$-Bereich bestimmen (Umkehrung
      $x{=}g^{-1}(y)$) und $f_Y=f_X(x)/|g'|$ einsetzen.}
  \rs{$y$-Grenzen aus den $x$-Grenzen über $y{=}g(x)$ umrechnen.}
  \rs{Kontrolle: $\int f_Y\,\mathrm dy=1$.}
\end{rezept}
\fallstrick{Steigung-Betrag $|g'|$ (Vorzeichen egal).\; $y$-Achse \emph{umrechnen},
  nicht $x$-Grenzen übernehmen.\;
  \textbf{Rückwärts}: Steigung so wählen, dass $Y$
  \emph{gleichverteilt} wird $\Rightarrow g'=\dfrac{f_X}{f_Y}$ mit $f_Y=$const.}
\miniexample{stückw.\ lineare Kennlinie \;/\; Kennlinie $g(x){=}\sqrt{x}$}{%
  $X{\sim}\mathcal U[-2,2]$ ($f_X{=}\tfrac14$), $g'{=}2\,(x{<}0),\,4\,(x{>}0)\Rightarrow
  f_Y{=}\tfrac{1/4}{2}{=}\tfrac18\,(-4{<}y{<}0),\;\tfrac{1/4}{4}{=}\tfrac1{16}\,(0{<}y{<}8)$.\\
  $g{=}\sqrt{x}\,(x{>}0):\;g'{=}\tfrac{1}{2\sqrt{x}}$, $x{=}y^2\Rightarrow f_Y(y)=f_X(y^2)\cdot 2y$.}
\bereichende

\bereich[cGrundl]{Gemeinsame Verteilungen \& Korrelation}
\begin{formelreihe}
  \formelblock{Verbund-CDF}{F_{XY}(x,y)
    =\int_{-\infty}^{x}\!\int_{-\infty}^{y}\!
    f_{XY}(\xi,\eta)\,\mathrm{d}\eta\,\mathrm{d}\xi} &
  \formelblock{Randdichte $X$}{f_X(x)
    =\int_{-\infty}^{\infty} f_{XY}(x,y)\,\mathrm{d}y} &
  \formelblock{Randdichte $Y$}{f_Y(y)
    =\int_{-\infty}^{\infty} f_{XY}(x,y)\,\mathrm{d}x} \\[0.3em]
  \formelblock{Unabhängigkeit}{f_{XY}(x,y)=f_X(x)\cdot f_Y(y)
    \;\Leftrightarrow\; P(A\cap B)=P(A)\cdot P(B)} &
  \formelblock{Kovarianz}{\mathrm{Cov}(X,Y)
    =E\bigl\{(X{-}\mu_X)(Y{-}\mu_Y)\bigr\}
    =E\{XY\}-\mu_X\mu_Y\\[0.1em]
    {>}0\text{: steigen zus.}\\[-0.05em]
    {<}0\text{: gegläufig}\quad
    {=}0\text{: unkorr.}} &
  \formelblock{Korrelationskoeff.}{\rho_{XY}
    =\dfrac{\mathrm{Cov}(X,Y)}{\sigma_X\cdot\sigma_Y}\in[-1,1]\\[0.15em]
    {+}1\text{: perfekt pos. linear}\\
    \phantom{+}0\text{: kein lin. Zsmhang}\\
    {-}1\text{: perfekt neg. linear}} \\[0.3em]
  \formelblock{Orthogonalität \& Unkorr.}{%
    \text{orthog.: }E\{X\cdot Y\}=0\\[0.1em]
    \text{unkorr.: }\mathrm{Cov}(X,Y)=0\\[0.1em]
    \text{unabh.}\Rightarrow\text{unkorr., umgekehrt \textbf{nicht}!}} &
  \formelblock{Linearität EW}{E\{aX\}=a\,E\{X\},\quad
    \\E\{X{+}Y\}=E\{X\}+E\{Y\}} &
  \beispielblock{Zwei unabh.\ Würfel $X,Y$}{%
    P(X{=}1){=}\tfrac16,\;P(X{=}2){=}\tfrac36,\;P(X{=}3){=}\tfrac26\\
    P(Y{=}1){=}\tfrac46,\;P(Y{=}2){=}\tfrac26\\
    P(X{=}1\cap Y{=}1)=\tfrac16\cdot\tfrac46=\tfrac{4}{36}} \\
\end{formelreihe}
\bereichende

\bereich[cGrundl]{Summe von ZV \& Zentraler Grenzwertsatz}
\begin{formelreihe}
  \formelblock{Summe $Z=X+Y$ (unabh.)}{f_Z(z)=(f_X * f_Y)(z)
    =\int_{-\infty}^{\infty}\!f_X(\xi)\,f_Y(z{-}\xi)\,\mathrm{d}\xi\\[0.15em]
    \mu_Z=\mu_X+\mu_Y\quad\sigma_Z^2=\sigma_X^2+\sigma_Y^2} &
  \formelblock{Zentraler Grenzwertsatz}{Y=\dfrac{1}{N}\sum_{i=1}^{N}X_i
    \;\xrightarrow{N\to\infty}\;\mathcal{N}(\mu_X,\tfrac{\sigma_X^2}{N})\\[0.15em]
    X_i\text{ unabh., gleichverteilt (beliebig)}} &
  \infoblock{Summe vieler unabh.\ ZV $\to$ Gaußverteilt,\\
    unabhängig von Einzelverteilung.\\
    Histogramm $\to$ Dichte für $N\to\infty$.} \\
\end{formelreihe}
\bereichende

\bereich[cGrundl]{REZEPT · Summe $Z{=}X{+}Y$ über Faltung}
\noindent\footnotesize
$\mu_Z=\mu_X+\mu_Y$,\;\; $\sigma_Z^2=\sigma_X^2+\sigma_Y^2$ (unabh.),\;\;
$f_Z=f_X*f_Y$.
\begin{rezept}
  \rs{Verstärkung \emph{vor} der Summe zuerst transformieren ($Y{=}aX$:
      Dichte um $a$ stauchen/spiegeln, s.\ Transformations-Rezept).}
  \rs{Diskret${+}$kontinuierlich: jede diskrete Masse $w_i$ bei $x_i$ erzeugt eine
      \emph{um $x_i$ verschobene, mit $w_i$ skalierte Kopie} von $f_Y$.}
  \rs{Kopien addieren; in Überlappungsbereichen Höhen aufsummieren.}
  \rs{Kontrolle: Gesamtfläche $=1$, $\mu_Z=\mu_X+\mu_Y$.}
\end{rezept}
\miniexample{diskret${+}$kontinuierlich \;/\; mit Verstärkung}{%
  $X\in\{0,1\}$, $P(0){=}\tfrac14$, $P(1){=}\tfrac34$, $Y\sim\mathcal U[0,3]$ ($f_Y{=}\tfrac13$):\;
  $f_Z=\tfrac14{\cdot}\tfrac13$ auf $[0,3]\;+\;\tfrac34{\cdot}\tfrac13$ auf $[1,4]$
  (Überlapp $[1,3]:\tfrac1{12}{+}\tfrac14{=}\tfrac13$).\\
  $Z{=}{-}2X{+}Y$, $Y$ zwei Diracs bei $\pm1$ (je $\tfrac12$) $\Rightarrow
  f_Z=\tfrac12 f_{-2X}(z{-}1)+\tfrac12 f_{-2X}(z{+}1)$ (zwei verschobene Kopien).}
\bereichende

\bereich[cGrundl]{Gesetz der großen Zahlen}
\begin{formelreihe}
  \infoblock{Relative Häufigkeit nähert sich bei $N\to\infty$\\
             der Wahrscheinlichkeit.\\
             Histogramm $\to$ Dichte für $N\to\infty$.} &
  \formelblock{Relative Häufigkeit}{h_X(x)=\dfrac{n(x)}{N}} &
  \formelblock{Konvergenz}{\lim_{N\to\infty}h_X(x)=P(X{=}x)} \\
\end{formelreihe}
\bereichende

% (Rezepte stehen jetzt direkt unter ihren Theorie-Blöcken)
